\documentclass[a4paper, norsk, 12pt]{article}   
% ===================================
% Language 
% ===================================
    \usepackage[norsk]{babel}
    \usepackage[utf8]{inputenc}

% ===================================
% Typography
% ===================================
    \usepackage[margin=1in,includeheadfoot]{geometry}% <- instead package fullpage  
    \usepackage{%
    			lmodern,%
    			fix-cm}% Better typography, e.g. kerning
    			
    \usepackage[activate={true,nocompatibility},
    			final,
    			tracking=true,
    			kerning=true,
    			spacing=true,
    			factor=1100,
    			stretch=10,
    			shrink=10]{microtype} \LoadMicrotypeFile{bch}
    \microtypecontext{spacing=nonfrench} 
    
% ===================================
% Figures
% ===================================
\usepackage{graphicx}
\usepackage{subcaption}
    
% ===================================
% Mathematics
% ===================================
    \usepackage{amsmath, amssymb, amsthm, mathtools}   % Mijøer og verktøy for å
    % skrive matematikk
    \usepackage{esdiff}
    \let\div\relax
    \DeclareMathOperator{\div}{div}
    % \DeclareMathOperator{\curl}{curl}
    
    % Øistein trenger flere egendefinerte kommandoer
    
    \DeclareMathOperator{\curl}{curl}
    \DeclareMathOperator{\grad}{grad}
    
    \newcommand{\vek}[1]{\mathbf{#1}}
    \newcommand{\unitvek}[1]{\vek{\hat{#1}}}
    \newcommand{\normal}{\mathbf{n}}
    
    \DeclarePairedDelimiter\abs{\lvert}{\rvert}%
    \DeclarePairedDelimiter\norm{\lVert}{\rVert}%

    \makeatletter
    \newcommand*{\rom}[1]{\expandafter\@slowromancap\romannumeral #1@}
    \newcommand*{\dd}{\mathop{}\!{\operator@font d}}
    \makeatother
    
    \newcommand{\R}{\mathbb{R}}
    \newcommand{\N}{\mathbb{N}}
    
    \newcommand{\I}{\vek{i}}
    \newcommand{\J}{\vek{j}}
    \newcommand{\K}{\vek{k}}
    \newcommand{\X}{\vek{x}}
    \newcommand{\Y}{\vek{y}}
    \newcommand{\A}{\vek{a}}
    \newcommand{\B}{\vek{b}}
    
    \newcommand{\dt}{\mathop{}\! \mathrm{d} t\mathop{}\! }
    \newcommand{\du}{\mathop{}\! \mathrm{d} u\mathop{}\! }
    \newcommand{\dv}{\mathop{}\! \mathrm{d} v\mathop{}\! }
    \newcommand{\ds}{\mathop{}\! \mathrm{d} s\mathop{}\! }
    \newcommand{\dx}{\mathop{}\! \mathrm{d} x\mathop{}\! }
    \newcommand{\dy}{\mathop{}\! \mathrm{d} y\mathop{}\! }
    \newcommand{\dA}{\mathop{}\! \mathrm{d} A\mathop{}\! }
    \newcommand{\dV}{\mathop{}\! \mathrm{d} V\mathop{}\! }
    \newcommand{\dS}{\mathop{}\! \mathrm{d} S\mathop{}\! }
    \newcommand{\dSS}{\mathop{}\! \mathrm{d} \vek{S}\mathop{}\! }

% ===================================
% Header and footer 
% ===================================
    \usepackage{lastpage} % number of last page 
    \usepackage{fancyhdr}
    
    \pagestyle{fancy}
    \fancyhf{}
    \renewcommand\footrule{\begin{minipage}{1\textwidth}
    \hrule width \hsize \kern 1mm \hrule width \hsize height 2pt  
    \end{minipage}\par}%
    
    \renewcommand\headrule{\begin{minipage}{1\textwidth}
    \hrule width \hsize height 2pt \kern 1mm \hrule width \hsize   
    \end{minipage}\par}%
    
    \cfoot{\today}
    \rfoot{Side \thepage\ av \pageref{LastPage}}
    \lhead{MA1103 Flerdimensjonal analyse, vår 2018}
    
    \fancypagestyle{plain}{% Redefine ``plain'' style for chapter boundaries  
      \fancyhf{} % clear all header and footer fields  
      \renewcommand\footrule{\begin{minipage}{1\textwidth}
      \hrule width \hsize \kern 1mm \hrule width \hsize height 2pt 
      \end{minipage}\par}%
      \renewcommand{\headrulewidth}{0pt}
      %
      \renewcommand{\headrule}{}% <- added
      %
      \fancyfoot[C]{\today} % except the center  
      \fancyfoot[R]{Side \thepage\ av \pageref{LastPage}} % except the center  
    }

% ===================================
% Problems 
% ===================================
    \usepackage{xparse}
    \usepackage{nameref}
    \newcounter{problem}
    \newcounter{problemMA1103}
    
    %\DeclareDocumentCommand{\oppgave}{o}{Oppgave~\stepcounter{problemMA1103}\arabic{problemMA1103}}

    % Defines a simple bold text, \oppgave[text] produces Oppgave 1 (text). 
    \NewDocumentCommand{\oppgave}{o}{%
      % <code>%
      \noindent
      % \textbf{\large Oppgave~\stepcounter{problemMA1103}\arabic{problemMA1103}}
      \stepcounter{problem\arabic{section}}\paragraph{Oppgave~\arabic{problem\arabic{section}}}
      \IfNoValueTF{#1}
        {}
        {(#1)}%
        % \medskip\\\noindent
        % <code>
        \hspace{0.1cm}
      }
      \usepackage{totcount}
      \regtotcounter{section}

      \newcommand{\seksjon}[1]{\section{#1}\label{sec:\arabic{section}}\newcounter{problem\arabic{section}}}

    
    % Changes the style for enumerate
    \usepackage{enumitem}
    \setlist[enumerate,1]{label={\alph*)}}
    \setlist[enumerate,2]{label={\roman*)}}

% ===================================
% Titlebox 
% ===================================
    \usepackage{tikz,pgfplots}   

    % Defines the NTNU logo
    \definecolor{mblue}{RGB}{0,72,161}
    \newcommand{\ntnulogo}{%
      \leavevmode
      \vtop{
        \hrule height 0pt
        \hbox{%
          \begin{tikzpicture}
          \fill[mblue, rounded corners=5mm] (0,2) rectangle (2,4);
          \fill[white] (0.5,2.5) rectangle (1.5,3.5);
          \fill[mblue] (1,3) circle (4mm);
          \end{tikzpicture}%
        }
      }
    }
    
    % Rename this for other subjects / year
    \newcommand{\course}{%
      \leavevmode
      \vtop{\LARGE
        \hrule height 0pt
        \ialign{\hfil##\cr
          MA1103\cr
          Flerdimensjonal analyse\cr
          Vår 2018\cr
        }%
      }%
    }
    
    % Creates the titlebox
    % \titlebox{Øving 1}{Frist ...} or \titlebox{Løsningsforslag - Øving 1}{}
    \newcommand{\titlebox}[2]{
    \thispagestyle{plain}
    \begin{center}
    \setlength{\fboxsep}{0pt}
    \fbox{%
      \begin{minipage}{\dimexpr\textwidth-2\fboxrule}
      \vspace{1pt}
      \hrule height 1pt
      \vspace{6pt}
      \mbox{}\quad\ntnulogo\hfill\course\quad\mbox{}\\[3pt]
      \mbox{}\quad
      \begin{tabular}{@{}l@{}}
      Norges teknisk-vitenskapelige\\
      universitet\\
      Institutt for matematiske fag
      \end{tabular}\hfill%
      \textbf{ #1. #2}\quad\mbox{}\\
      \vspace{3pt}
      \hrule height1pt
      \vspace{1pt}
      \end{minipage}%
    }
    \end{center}
    }

    
\begin{document}


\seksjon{Grenseverdier}
  
\oppgave[V2017, Oppgave 1] La
%
\begin{equation*}
  f(x, y) :=
  \begin{cases}
    \cfrac{xy^2}{x^2 + y^2}, & \text{når} \quad (x,y) \neq (0,0) \\
    0, & \text{når} \quad (x,y) = (0,0)
  \end{cases}
\end{equation*}

\begin{enumerate}
  \item Vis at
    \begin{enumerate}
      \item $f$ er kontinuerlig i $(x,y) = (0,0)$,
      \item $\diffp{f}{x}(0,0)$ og $\diffp{f}{y}(0,0)$ eksisterer, men $f$ er
        ikke deriverbar i $(0,0)$.
    \end{enumerate}
    % 
  \item La $g(t) = (at, bt)$ med konstanter $a$ og $b$ ulik $0$ og verifiser at
    %
    \begin{equation*}
      (f \circ g)'(0) = \frac{ab^2}{a^2 + b^2},
      \qquad \text{men} \qquad 
      \nabla f(0,0) \cdot g'(0) = 0\,.
    \end{equation*}
    %
    Forklar at dette ikke er i strid med kjerneregelen.
\end{enumerate}

\oppgave[K2016, Oppgave 1]

\begin{enumerate}
  \item Vis at $\displaystyle \lim_{(x,y) \to (0,0)} \frac{x^2y}{x^2 + y^2} =
    0$.
  \item La $f(w,x,y,z) = w - x^2y^3z$. Beregn grenseverdien:
    %
    \begin{equation*}
      \lim_{h \to 0} \frac{f(5,2,1+h,-1) - f(5,2,1,-1)}{h}.
    \end{equation*}
\end{enumerate}


\oppgave[V2016, Oppgave 1] La
%
\begin{equation*}
  f(x, y) :=
  \begin{cases}
    \cfrac{2xy}{x^2 + y^2}, & \text{når} \quad (x,y) \neq (0,0), \\
    0, & \text{når} \quad (x,y) = (0,0)
  \end{cases}\,.
\end{equation*}
%
Forklar hvorfor $f$ ikke er deriverbar i $(0,0)$, men de partiellderiverte
$\diffp{f}{x}(0,0)$ og $\diffp{f}{y}(0,0)$ eksisterer.

\oppgave[K2015, Oppgave 1] La
%
\begin{equation*}
  f(x, y) :=
  \begin{cases}
    \cfrac{x^2y^4}{x^2 + y^2}, & \text{når} \quad (x,y) \neq (0,0), \\
    0, & \text{når} \quad (x,y) = (0,0)
  \end{cases}\,.
\end{equation*}
%
\begin{enumerate}
  \item Vis at $\diffp{f}{x}(0,0)$ og $\diffp{f}{y}(0,0)$ eksisterer
  \item Vis at $f$ ikke er deriverbar i $(0,0)$ ved å vise at $f$ ikke er
    kontinuerlig i $(0,0)$.
\end{enumerate}

\oppgave[K2014, Oppgave 3] Begrunn at funksjonen
%
\begin{equation*}
  f(x, y) :=
  \begin{cases}
    \cfrac{\sin \sqrt{x^2+y^2\,}\,}{\sqrt{x^2 + y^2\,}\,}, & \text{når} \quad (x,y) \neq (0,0), \\
    0, & \text{når} \quad (x,y) = (0,0)
  \end{cases}\,,
\end{equation*}
%
er kontinuerlig.

\oppgave[K2013, Oppgave 3] Vis at en av grensene under eksisterer mens den andre
ikke eksisterer
%
\begin{equation*}
  \lim_{(x,y)\to(0,0)} \frac{xy}{\sqrt{x^2+y^2}}
  \qquad
  \lim_{(x,y)\to(0,0)} \frac{xy}{x^2 + y^2}.
\end{equation*}


\oppgave[K2012, Oppgave 5] Funksjonen $f$ er gitt ved
%
\begin{equation*}
  f(x, y) :=
  \begin{cases}
    \cfrac{2xy}{\sqrt{x^2 + y^2\,}\,}, & \text{når} \quad (x,y) \neq (0,0), \\
    0, & \text{når} \quad (x,y) = (0,0)
  \end{cases}\,,
\end{equation*}
%
\begin{enumerate}
  \item Vis at $f$ er kontinuerlig i origo.
  \item En kan vise at $f$ ikke er kontinuerlig deriverbar i origo. Vis likevel
    at den retningsderiverte
    %
    \begin{equation*}
      D_{\vek{u}}f(0,0) = \lim_{t \to 0^+} \frac{f(t u_1, t u_2) - f(0,0)}{t}
    \end{equation*}
    %
    eksisterer for alle enhetsvektorer $\vek{u} = u_1 \I + u_2 \J$.
\end{enumerate}

\oppgave[V2010, Oppgave 3] Gitt funksjonen
%
\begin{equation*}
  f(x, y) :=
  \begin{cases}
    \cfrac{Ax^3 + By^3 - xy}{2x^2 + 2y^2} & \text{når} \quad (x,y) \neq (0,0), \\
    0, & \text{når} \quad (x,y) = (0,0)
  \end{cases}\,.
\end{equation*}
%
\begin{enumerate}
\item Bestem konstanter $A$ og $B$ slik at $\diffp{f}{x}(0,0) = 3$ og
$\diffp{f}{y}(0,0) = 0$. \medskip

\item La nå $A = 6$ og $B = 0$. Er $f$ kontinuerlig i origo for dette valget av konstanter?
\end{enumerate}

%%% Local Variables:
%%% mode: latex
%%% TeX-master: "main"
%%% End:

\seksjon{Gradient}

\oppgave[V2015, Oppgave 2]

Funksjonen $f \colon \R^2 \to \R$ er deriverbar. Vi vet også at den
retningsderiverte i $(1, 0)$ langs positiv $x$-akse er $5$, og at den
retningsderiverte i $(1, 0)$ langs linja $y = x - 1$ i retning av positiv $y$,
er $- 2$. Hva er gradienten til $f$ i $(1, 0)$?

\oppgave[K2014, Oppgave 5]

Anta at et fjell har form som en elliptisk paraboloide $z = c - ax^2 - by^2$,
der $a$, $b$ og $c$ er positive konstanter, $x$ og $y$ er øst–vest og nord–sør
koordinater på kartet, mens $z$ er høyden over havet.
%
\begin{enumerate}
  \item I hvilken retning stiger høyden mest i punktet $(1, 1)$?
  \item Hvis en klinkekule slippes i $(1, 1, c - a - b)$, i hvilken retning vil den begynne
    å trille?
\end{enumerate}

\oppgave[V2014, Oppgave 2]

La $T(x,y,z) = e^{x + 2y + 3z}$ være temperaturen i et romlig område om origo.
\medskip

I hvilke retninger fra $(0,0,0)$ vokser og avtar temperaturen mest? \medskip

Hva er den retningsderiverte i disse retningene?

\oppgave[V2012, Oppgave 2] Finn den retningsderiverte av funksjonen $f(x,y,z) =
e^{-x^2}y - \log(1 + e^z)$ i punktet $(1,1,0)$ i retningen fra $(1,1,0)$ til $(-1,2,1)$.
\documentclass[tikz, border = 2mm]{standalone}
% \usepackage{kompendium-figurer}
\usepackage{pgfplots}
\usepgfplotslibrary{fillbetween}

    \pgfkeys{/pgfplots/Delta Style/.style={
		    % scale only axis,
		    grid=major,
		    axis equal,
		    grid style={dashed, gray!30}, %Uncomment these lines for no grid
		    axis lines=middle,
		    inner axis line style={=stealth}, %Arrow type
        ultra thick,
		    xlabel={\large $x$},
		    ylabel={\large $y$},
        cycle list = {black,black!70,black!40,black!10} %Plot colors cycle in grayscale
      }}

    \pgfplotsset{
        % use this `compat' level or higher to use the LUA backend for calculation
        % (--> speed improvement)
        compat=1.12,
        % declare your function here ...
        /pgf/declare function={
            f(\x) = 2*\x+1;
        },
        /pgf/declare function={
            y(\y) = 0.5*(\y-1);
        },
    }

\begin{document}
\begin{tikzpicture}
  \def\e{1}
  \def\ea{3-\e}
  \def\eb{3+\e}
  
  \def\d{0.75}
  \def\da{1-\d}
  \def\db{1+\d}
  
  \begin{axis}[
    Delta Style,
    legend pos =north west,
    ytick={-1,0,...,8},
    xtick={-1,0,...,3},
    ymin=-0.5,
    ymax=5.5,
    xmin=-1.5,
    xmax=2.5,
		]
    % Filling
    \addplot[name path=plot,ultra thick,samples=100,domain= -10:10] {f(x)};
    \addlegendentry{$y=2x+1$}
    \path[name path=yaxis](current axis.below origin)-- (current axis.above origin);

    \addplot[red!50, fill opacity=0.5] fill between[
    of=yaxis and plot,reverse=false,
    soft clip={domain y = \ea:\eb},
    ];

    \addplot[domain=\da:\db,blue,dashed, fill=blue!50, fill opacity=0.5] {f(x)}\closedcycle;
  
        \addplot[color=gray,fill=white,only marks,mark size = 4pt] coordinates {(1,3)};

        \addplot[dashed,thick,red, domain=-2.5:y(\ea)] {\ea};
        \addplot[dashed,thick,red, domain=-2.5:y(\eb)] {\eb};

        \draw[ultra thick, red, stealth-stealth] %
        (axis cs: -1, \ea) -- node[left] { $\varepsilon$ } (axis cs:-1, \eb);
        \draw[ultra thick, blue, stealth-stealth] %
        (axis cs: \da, 1) -- node[below] { $\delta$ } (axis cs: \db, 1);
		\end{axis}
	\end{tikzpicture}
\end{document}
\newpage


\seksjon{Kjerneregelen}

\oppgave[V2016, Oppgave 2] Anta at en flue beveger seg langs en kurve $C$ i
$\R^3$ slik at posisjonen til flue er gitt ved $\vek{c}(t) = t \I +
\frac{2}{3}\sqrt{2} t^{3/2} \J + \frac{1}{2}t^2 \K$ ved tiden $t
\subset [0, \infty)$.
    
\begin{enumerate}
    \stepcounter{enumi}
  \item La $T(x,y,z) = x^2 + xz + y$ være temperaturen i punktet $(x,y,z) \in
    \R^3$. Finn $\diff{T}{t} \vek{c}(t)$ altså temperaturendringen flua vil
    oppleve ved tiden $t = 1$.
\end{enumerate}

\oppgave[V2016, Oppgave 6] Med definisjonen $\nabla^2 := \nabla \cdot \nabla =
\diffp[2]{}{x} + \diffp[2]{}{y}$, kalles en $C^2$-funksjon $f \colon \R^2 \to
\R$ som tilfredstiller ligningen
%
\begin{equation*}
  \nabla^2 f = 0
\end{equation*}
%
for \emph{harmonisk}. Ved at for et vilkårlig reelt tall $k$ er funksjonen
$f(x,y) = e^{kx}(\cos ky)$ harmonisk.


\oppgave[V2015, Oppgave 4] Vis at hvis akselerasjonen $\vek{a}(t)$ til en
punktmasse alltid er perpendikulær på hastigheten $\vek{v}(t)$, så er farten
$\norm{\vek{v}}$ konstant. (Hint: $\norm{\vek{v}}^2 = \vek{v} \cdot \vek{v}$.


\oppgave[K2014, Oppgave 2] Gitt $z = f(x, y)$ der $f$ er en deriverbar funksjon
og $x = u^2 - v^2$, $y = v^2 - u^2$. Vis at
%
\begin{equation*}
  u \diffp{z}{v} + v \diffp{z}{u} = 0\,.
\end{equation*}

\oppgave[V2014, Oppgave 1] Gitt $z = f(x,y)$ der $f$ er en deriverbar funksjon,
$x = u + v$ og $y = u - v$. Finn $\diffp{z}{u}$, $\diffp{z}{v}$ og vis at
%
\begin{equation*}
  \diffp{z}{u} \cdot \diffp{z}{v}
  = \left( \diffp{z}{x} \right)^2 - \left( \diffp{z}{y} \right)^2.
\end{equation*}

\oppgave[V2013, Oppgave 1] Fra fysiske lover kan en se at om $K$ er et homogent
legeme i $\R^3$, så må temperaturen $T = T(x,y,z)$ i $K$ være en løsning til
\emph{varmelikningen}
%
\begin{equation*}
  k \left( \diff[2]{T}{x} + \diff[2]{T}{y} + \diff[2]{T}{z} \right)
  = \diffp{T}{t}
\end{equation*}
%
der $(x,y,z)$ er posisjonen i legemet, $t$ er tiden, og $k$ er en
materialkonstant. \medskip

Vis at $T(x,y,z,t) = 2x - y + z$ er en løsning til varmelikningen.

\oppgave[V2012, Oppgave 4] Vi sier at $f(x,y)$ er harmonisk hvis $f$ er to
ganger kontinuerlig deriverbar og
%
\begin{equation*}
  \diffp[2]{f}{x} + \diffp[2]{f}{y} = 0.
\end{equation*}
%
Vis at hvis $g$ er harmonisk så er også $f(x,y) = g(x^2 - y^2, 2xy)$ harmonisk.

% \oppgave[V2013, Oppgave 4] Vis sier at en funksjonen er $C^2$ om alle
% andreordens partiellderiverte er kontinuerlige. Vis at om $f$ er en $C^2$
% funksjon, så er $\curl \nabla f = \mathbf{0}$. Si klart fra hvordan du bruker at
% funksjonen er $C^2$. 


\seksjon{Taylor-approksimasjon}

\oppgave[K2016, Oppgave 3] La $f(x, y) = x^4 + x^3 + y^2 + xy$. Finn andreordens
Taylor-approksimasjon av $f$ i punktet $(x_0, y_0) = (1, 2)$.
  

\seksjon{Tangenter}

\oppgave[K2016, Oppgave 2] La $C$ være en kurve som er parametrisert ved
%
\begin{equation*}
  \vek{c}(t) = (t^2, 2t, 4 - t), \qquad -2 \leq t \leq 2.
\end{equation*}
%
\begin{enumerate}
  \item Vis at punktet $(0,0,4)$ ligger på kurven og finn en tangentvektor til
    kurven i dette punktet som har lengde $1$.
  \item Finn en ligning for tangentplanet som er ortogonalt til kurven i punktet $(0,0,4)$.
\end{enumerate}

\oppgave[K2015, Oppgave 3] Finn ligningen for tangentplanet til flaten $x^2 +
y^2 - e^{xz} - \sin y = 0$ i punktet $(1,0,0)$ og bruk denne til å finne en
tilnærmet verdi for $x$ i det punktet på flaten som ligger i nærheten av
$(1,0,0)$ med $y = z = 1/10$.

\oppgave[V2015, Oppgave 1] Finn en ligning for tangentplanet til ellipsoiden
%
\begin{equation*}
  2x^2 + 3y^2 + z^2 = 9
\end{equation*}
%
i punktet $(1, -1, 2)$.

\oppgave[K2014, Oppgave 1] Finn ligningen for tangentplanet til flaten \\ $x^2 +
z^2 - e^{xy} - \sin z = 0$ i punktet $(1, 0, 0)$.

\oppgave[K2013, Oppgave 1] Finn en ligning for tangentplanet til flaten \\ $z^2 =
2x^2 - y^2$ i punktet $(1, -1, 1)$.

\oppgave[V2013, Oppgave 3] En kurve $C$ er parametrisert med
%
\begin{equation*}
  \vek{c}(t) = t^2\I + 2t \J + (4 - t) \K, \qquad -2 \leq t \leq 2\,.
\end{equation*}
%
\begin{enumerate}
  \item Vis at punktet $(1,2,3)$ ligger på kurven $C$, og finn en tangentvektor
    til kurven i dette punktet, med lengde $1$.
  \item Finn en likning for planet som står normalt på kurven i punktet $(1,2,3)$.
\end{enumerate}

\oppgave[V2012, Oppgave 1] Finn en parameterfremstilling for den rette linja
gjennom $P = (1,2,3)$ parallell med vektoren $2\I + \J - \K$.

\seksjon{Kritiske punkter}
  
\oppgave[V2017, Oppgave 3]

\begin{enumerate}
  \item Finn og klassifiser alle kritiske punkt til $f(x, y) = x^3 - y^2 - 3x - 6y - 1$.
  \item Finn den største og minste verdien til $f(x, y) = x^2 - xy + y^2$ på
    kvartsirkelen $x^2 + y^2 = 1$, $x, y \geq 0$.
\end{enumerate}


\oppgave[K2016, Oppgave 4] La $f \colon \R^2 \to \R$ være gitt ved $f(x,y) =
e^{xy}$.

\begin{enumerate}
  \item Finn og klassifiser alle kritiske punkter til $f$.
  \item Hvor oppnår $f$ sitt maksimum og minimum når $2x^2 + y^2 = 1$?
\end{enumerate}


\oppgave[V2016, Oppgave 3]

\begin{enumerate}
  \item Finn og klassifiser alle kritiske punkter av $f(x, y) = 2x^2 + 4xy +
    y^4$.
  \item Hvor oppnår $f(x,y) = (x - y)^2$ sitt maksimum og minimum når $x^2 + y^2
    = 2$?
\end{enumerate}


\oppgave[K2015, Oppgave 4] Funksjonen $f$ er gitt ved $f(x, y) = 2x^2 - x^4 + y^2$.

\begin{enumerate}
  \item Finn alle kritiske punkter til $f$, og avgjør om disse er lokale
    maksima, minima eller saddelpunkter.
  \item Finn største og minste verdi for $f$ på kurven $x^4 + y^2 = 4$. 
\end{enumerate}


\oppgave[V2015, Oppgave 3] Bestem den største og minste verdien som funksjonen
$f(x,y) = xy$ oppnår i området
%
\begin{equation*}
  D = \{ (x,y) \in \R^2 \mid x^2 + xy + y^2 \leq 3 \}.
\end{equation*}


\oppgave[K2014, Oppgave 6] La $f(x,y) = 3x^2 y - y^3 + x^4$.

\begin{enumerate}
  \item Finn alle kritiske punkt til $f$ og klassifiser disse som lokale
    maksimum, lokale minimum eller saddelpunkt.
  \item Har $f$ et globalt (absolutt) maksimum? Et globalt (absolutt) minimum?
\end{enumerate}


\oppgave[K2013, Oppgave 3] Finn og klassifiser alle de kritiske punktene til
%
\begin{equation*}
  f(x,y) = xe^{-x^2-y^2}.
\end{equation*}


\oppgave[V2013, Oppgave 4] Vi lar funksjonen $f(x,y) = x^3 - xy^2 - 2x^2 + y^2$
ha hele planet som definisjonsområdet.
%
\begin{enumerate}
  \item Finn alle kritiske punkt til $f$ og klassifiser disse som lokale
    maksimum, lokale minimum eller saddelpunkt.
  \item Har $f$ globale maksimum eller minimum? Finn i så fall disse. 
\end{enumerate}

\oppgave[K2012, Oppgave 2] Finn største og minste verdi til funksjonen \\ $f(x,y,z)
= z$ langs skjæringskurven mellom flatene $x^2 + 2y^2 = 1$ og $z = x - 4y$.


\oppgave[V2012, Oppgave 3] Finn maximumsverdien av funksjonen
%
\begin{equation*}
  f(x,y) = 2x + y - x^2 - 2y^2 + 3.
\end{equation*}


\oppgave[V2012, Oppgave 6] Finn minimumsverdien av funksjonen $f(x,y) = xy$ på
ellipsen $g(x,y) = x^2 + 2y^2 $.




%%% Local Variables:
%%% mode: latex
%%% TeX-master: "main"
%%% End:

\newpage
\seksjon{Optimering}

\oppgave[V2017, Oppgave 2] La $t$ betegne tiden. Banen til romskip $A$ er gitt ved
%
\begin{equation*}
  \vek{r}_A(t) = (t, t^2, t^2 - t)
\end{equation*}
%
og banen til romskip $B$ er gitt ved
%
\begin{equation*}
  \vek{r}_B(t) = (8 - t, t^2 - 4, t^2 - 2t).
\end{equation*}
%
Ved hvilken tid $t \geq 0$ er det minst avstand mellom skipene?

\oppgave[K2014, Oppgave 7] Finn punktene på kuleflaten $x^2 + y^2 + z^2 = 4$ som
er nærmeste og lengst fra punktet $(2,2,1)$.

\oppgave[K2013, Oppgave 5] Finn det punktet på planet $3x - 2y + z = 7$ som
ligger nærmest origo.

\oppgave[V2013, Oppgave 2] I USA definerer postverket ``størrelsen'' (``size'')
til en pakke som summen av lengde (length) og ``girth'', der ``girth'' er
omkretsen normal på lengden. ``Størrelsen''/``size'' kan ikke være mer enn $130$
tommer. \medskip

Hva er største volum (i kubikktommer) en rektangulær pakke kan ha?


%%% Local Variables:
%%% mode: latex
%%% TeX-master: "main"
%%% End:


\seksjon{Linjeintegral}

\oppgave[K2016, Oppgave 5] La $\vek{F}(x, y) = (0, x)$.

\begin{enumerate}
  \item Hva er verdien av integralet
    %
    $ \displaystyle
      \int_C \vek{F} \cdot \dd\vek{r},
    $
    %
    der $C$ er en sirkel med radius $a>0$?
\end{enumerate}

\oppgave[V2016, Oppgave 2] Anta at en flue beveger seg langs en kurve $C$ i
$\R^3$ slik at posisjonen til flue er gitt ved $\vek{c}(t) = t \I +
\frac{2}{3}\sqrt{2} t^{3/2} \J + \frac{1}{2}t^2 \K$ ved tiden $t
\subset [0, \infty)$.
    
\begin{enumerate}
  \item Dersom flue starter med å fly ved $t = 0$ og farten er gitt i meter per
    sekund, hvor mange meter tilbakelegger flue i løpet av $10$ sekunder?
\end{enumerate}

\oppgave[K2015, Oppgave 2] Finn buelengden til kurven med parameterfremstilling
%
\begin{equation*}
  \vek{r}(t) = (\cos t + t \sin t)\I + (\sin t - t \cos t)\J + t^2 \K, \qquad \text{for} \ 0 \leq t \leq 2\pi.
\end{equation*}


\seksjon{Konservative vektorfelt}

\oppgave[V2016, Oppgave 5] La $\vek{F}(x,y,z) = 2xy\I + x^2\J + z\K$. Beregn
linjeintegralet
%
\begin{equation*}
  \int_C \vek{F} \cdot \dd \vek{r},
\end{equation*}
%
\begin{enumerate}
  \item $C$ er kurven parametrisert ved $\vek{r} \colon [0, 2\pi] \to \R^3$,
    $\vek{r} = (\cos t) \I + (\sin 2t)\J + (\sin t)^2\K$.
  \item $C$ er en vilkårlig glatt kurve med startpunkt $(0,0,0)$ og endepunkt
    $(1, -2, \sqrt{2})$.
\end{enumerate}


\oppgave[K2015, Oppgave 6] La $\vek{F}(x,y,z) = 3x^2y\I + (x^3 + y^3)\J$.
Verifiser at $\curl \vek{F} = 0$ og finn en funksjon $f$ slik at $\vek{F} =
\grad f$.

\oppgave[V2015, Oppgave 8] La $\vek{F}(x,y,z)$ være et vektorfelt i $\R^3$ gitt ved
$\vek{F}(x,y,z) = yz \I + xz \J + xy\K$. Bestem
%
\begin{equation*}
  \int_C \vek{F} \cdot \dd\vek{r}
\end{equation*}
%
når $C$ er kurven med parametrisering $\vek{r}(t) = (\cos t) \I + (\sin t)J +
t\K$, $0 \leq t \leq \pi/4$.


\oppgave[V2014, Oppgave 6] Gitt vektorfeltet $\vek{F}(x,y,z) = x^2I + xy\J +
(2x^2 - z^2)K$.
%
\begin{enumerate}
  \item Finn $\curl \vek{F}$ og $\div \vek{F}$.
  \item Grunngi at $\vek{F}$ hverken er gradienten til en $C^2$-funksjon, eller
    curl til et $C^2$-vektorfelt.
\end{enumerate}

\oppgave[K2013, Oppgave 2] Vektorfeltet $\vek{F}(x,y,z) = 2xy\I + x^2\J + \K$ er
definert for alle $(x,y,z) \in \mathbb{R}^3$ i rommet.

\begin{enumerate}
  \item Vis at $\vek{F}$ er et konservativt vektorfelt.
  \item La $C$ være kurven i rommet med parametrisering $\vek{c}(t) = t^3\I +
    2\J + (1+t^2)\K$, $0 \leq t \leq 1$. Finn
    %
    \begin{equation*}
      \int_C \vek{F} \cdot \dd\vek{r}
    \end{equation*}
\end{enumerate}

\oppgave[V2013, Oppgave 4] Vis sier at en funksjonen er $C^2$ om alle
andreordens partiellderiverte er kontinuerlige. Vis at om $f$ er en $C^2$
funksjon, så er $\curl \nabla f = \mathbf{0}$. Si klart fra hvordan du bruker at
funksjonen er $C^2$. Dette viser med andre ord at et konservativt vektorfelt
$\vek{F} = \nabla f$ er \emph{rotasjonsfritt}.


\oppgave[K2012, Oppgave 1] La $\vek{F}(x,y,z)$ være definert for alle $(x,y,z)
\in \R^3$ ved
% 
\begin{equation*}
  \vek{F}(x,y,z) = 2xyz \I + z(x^2 - 4yz^2)\J + y(x^2 - 6yz^2)\K.
\end{equation*}
%
\begin{enumerate}
  \item Vis at $\vek{F}$ er et konservativt felt.
  \item Regn ut
    %
    \begin{equation*}
      \int_C \vek{F} \cdot \dd\vek{r}
    \end{equation*}
    %
    der $C$ er kurven i rommet gitt ved $\vek{r}(t) = (1+t^3)\I + t^2\J + \K$,
    $0 \leq t \leq 1$.
\end{enumerate}

\oppgave[V2012, Oppgave 7] Finn arbeidet utført av kraften $\vek{F}(x,y,z) =
yz\I + xz\J + xy\K$ på en partikkel som beveger seg på kurven $\vek{r}(t) =
e^t\I + (\sin t)\J + t^3\K$ i tidsintervalet $t \in [0, 1]$.

%%% Local Variables:
%%% mode: latex
%%% TeX-master: "main"
%%% End:


\seksjon{Dobbelintegraler}

\oppgave[V2017, Oppgave 4]

\begin{enumerate}
  \item Beregn dobbelintegralet $\displaystyle \iint_R \frac{x}{x^2 + y^2} \dd (x,y)$, der $R$
    er gitt ved $x^2 + y^2 \leq 1$, $x \geq 1$ og $y \geq 1$.
\end{enumerate}

\oppgave[V2016, Oppgave 4] La $D$ være området i første kvadrant ($x \geq 0, y
\geq 0$) begrenset av $x = 0$, $y = 0$ og $y = \sqrt{1 - x^2}$. Beregn
integralet av $f(x,y) = x^2 + y^2$ over $D$.


\oppgave[V2016, Oppgave 4] Beregn dobbelintegralet
%
\begin{equation*}
  \iint_D (2x + y^2) \dd (x,y)
\end{equation*}
%
når $D$ er alle punkter i første kvadrant som ligger inni sirkelskiven $x^2 +
y^2 \leq 4$, men utenfor enhetskvadratet $0 \leq x \leq 1$, $0 \leq y \leq 1$.

\oppgave[V2015, Oppgave 5] Gitt $g \colon D \to \R$ hvor
%
\begin{equation*}
  D = \{ (x,y) \in \R^2 \mid 1 \leq x \leq 9 \ \text{og} -3\sqrt{x} \leq y \leq 3 \sqrt{x} \}
\end{equation*}
%
og $g(x, y) = 1 + y^2/2x$.

\begin{enumerate}
  \item Finn volumet av området
    %
      $ \displaystyle T = \{ (x,y,z) \in \R^3 \mid (x,y) \in D \ \text{og} \ 0 \leq z \leq g(x,y)\}$.
\end{enumerate}


\oppgave[K2014, Oppgave 4] Bruk dobbelintegral til å finne arealet omsluttet av
kurven $r = 1 + \sin \theta$ (i polarkoordinater).


\oppgave[K2012, Oppgave 3] La $R$ være området i $xy$-planet som er avgrenset av
de to kurvene
%
\begin{equation*}
  3x^2 - y^2 = 3 \quad \text{for} \ x \geq 1
  \qquad \text{og} \qquad
  x + y^2 = 11.
\end{equation*}
%
Skissèr de to kurvene og beregn integralet $\displaystyle \iint_R x \dA$.

\oppgave[V2012, Oppgave 8] Finn integralet $\displaystyle \iint_D (x^2 + y^2)^7
\dx \dy$ hvor $D$ er sirkelen $x^2 + y^2 \leq 4$. 

%%% Local Variables:
%%% mode: latex
%%% TeX-master: "main"
%%% End:

\seksjon{Integrasjonsrekkefølge}

\oppgave[V2017, Oppgave 3]

\begin{enumerate}
    \setcounter{enumi}{2}
  \item Skisser integrasjonsområdet for
    %
    \begin{equation*}
      \int_0^1 \int_0^x f(x, y) \dy \dx\,,
    \end{equation*}
    %
    og bytt integrasjonsrekkefølgen til $\dx \dy$.
\end{enumerate}

\oppgave[K2016, Oppgave 5] Beregn dobbelintegralet
%
\begin{equation*}
  \int_0^1 \int_{3y}^3 e^{x^2} \dx \dy\,.
\end{equation*}

\oppgave[V2014, Oppgave 4] Regn ut integralet
%
\begin{equation*}
  \int_0^2 \int_x^2 e^{y^2} \dy \dx\,,
\end{equation*}
%
ved å bytte om integrasjonsrekkefølgen.

%%% Local Variables:
%%% mode: latex
%%% TeX-master: "main"
%%% End:

\newpage

\seksjon{Greens teorem}


\oppgave[V2017, Oppgave 6]

\begin{enumerate}
    \stepcounter{enumi}
  \item Bruk Greens's teorem og et passende vektorfelt $\vek{F}$ til å beregne
    arealet av den elliptiske skiven begrenset av kurven $C$ som er
    parametrisert ved
    %
    \begin{equation*}
      \vek{c}(t) = 3(\cos t + \sin t) \I + 2(\sin t - \cos t) \J, \qquad 0 \leq t < 2\pi.
    \end{equation*}
\end{enumerate}


\oppgave[V2014, Oppgave 5] Regn ut integralet
%
\begin{equation*}
  \int_C \bigl(5y - e^{\sin x}\bigr) \dx + \bigl(10x - \sin(y^3 + 8y)\bigr)\dy
\end{equation*}
%
der $C$ er en sirkel med radius $2$ og sentrum i $(a, b)$. Spesifiser om du
regner integralet med eller mot klokka.


\oppgave[K2013, Oppgave 6] La $C$ være randen til trekanten i $xy$-planet med
hjørner i $(0,0)$, $(0,1)$ og $(1,1)$. Finn
%
\begin{equation*}
  \int_C e^{y^2-x^2}(\cos 2xy)\dx + e^{y^2-x^2}(\sin 2xy) \dy\,.
\end{equation*}

\oppgave[V2012, Oppgave 9] Finn integralet
%
\begin{equation*}
  \int_C \bigl( e^{x^2} + y \bigr) \dx + \bigl( 2x - e^{y^2} \bigr) \dy
\end{equation*}
%
lans sirkelen $x^2 + y^2 = 4$. Spesifiser om du regner ut integralet med eller
mot klokken. 
\seksjon{Trippelintegral}


\oppgave[V2017, Oppgave 4]

\begin{enumerate}
    \stepcounter{enumi}
  \item Beregn trippelintegralet $\displaystyle \iiint_D e^{(x^2 + y^2 +
    z^{2})^{3/2}}\dd(x,y,z)$, der $D$ er gitt ved \\ $1 \leq x^2 + y^2 + z^2
    \leq 4$ og $z \geq 0$.
\end{enumerate}


\oppgave[K2013, Oppgave 7] La $T$ være legemet gitt ved $x \geq 0,$, $y \geq 0$
og $0 \leq z \leq 1 - \sqrt{x^2 + y^2}$. La $S$ være overflaten til $T$
orientert med enhetsnormal som peker ut av legemet $T$.

\begin{enumerate}
  \item Finn $\displaystyle \iiint_T z\bigl( x^2 + y^2 \bigr) \dV$.
\end{enumerate}


\oppgave[V2013, Oppgave 6] I denne oppgaven er

\begin{itemize}
  \item $S$ sfæren (kuleoverflaten) med sentrum i $(0,0,3)$ og radius $5$,
    orientert slik at normalvektoren peker ut av kula.
  \item $S^+$ den delen av $S$ som ligger over $xy$-planet, med samme
    orientering som $S$.
  \item $T$ området avgrenset av $S^+$ og $xy$-planet.
\end{itemize}

\begin{enumerate}
  \item Finn volumet til området $T$.
\end{enumerate}

\oppgave[K2012, Oppgave 6] Legemet $T$ er avgrenset av paraboloiden $z = 4x^2 +
4y^2$ og planet $z = 4$.
%
\begin{enumerate}
  \item Finn volumet til $T$.
\end{enumerate}

%%% Local Variables:
%%% mode: latex
%%% TeX-master: "main"
%%% End:


\seksjon{Divergensteoremet}

\oppgave[V2017, Oppgave 5]

\begin{enumerate}
  \item La $T$ være flata gitt ved $\{(x,y,z) \in \R^3 \mid x^2 + y^2 = 4
    \ \text{og} \ 0 \leq z \leq 1\}$ og $\vek{F}$ være vektorfeltet
    %
    \begin{equation*}
      \vek{F} = (x + \sin(z^2y), x^2 + z, 1 - z)
    \end{equation*}
    %
    Regn ut $\div \vek{F}$ og finn
    %
    \begin{equation*}
      \iint_T \vek{F} \cdot \vek{n} \dS 
    \end{equation*}
    %
    der normalvektoren $\vek{n}$ peker ut av $T$.
\end{enumerate}

\oppgave[K2016, Oppgave 7]

La $\vek{F}(x,y,z) = (bxy, bx^2y, (x^2 + y^2)z^2)$ og $W \subset \R^3$ være
sylinderen
%
\begin{equation*}
  W := \{  (x,y,z) \in \R^3 \mid x^2 + y^2 \leq a^2, z \in [0, b]\}
\end{equation*}
%
der $a, b > 0$. Bruk Gauss' teorem til å finne $\displaystyle \iint_{\partial W} \vek{F} \cdot
\dSS$.

\oppgave[V2016, Oppgave 7]

La $f \colon \R^3 \to R$ være en skalar $C^2$-funksjon og $\vek{F} \colon \R^3
\to R^3$ et $C^1$-vektorfelt. La $W \subset \R^3$ være et området der Gauss'
(divergensteoremet) gjelder og $\partial W$ overflata til $W$.

\begin{enumerate}
    \item Vis at $\div(f \vek{F}) = f \div \vek{F} + \nabla f \cdot \vek{F}$ og derfor
    %
    \begin{equation*}
      \iint_{\partial W} f \vek{F} \cdot \vek{n} \dS
      = \iint_{W} f \div \vek{F} \dV
      + \iint_{W} \nabla f \cdot \vek{F} \dV 
    \end{equation*}
    %
  \item Anta at
    %
    \begin{align*}
      \nabla^2 & f(x,y,z) = 0 \qquad \text{når} \qquad (x,y,z) \in W  \\
               & f(x,y,z) = 0 \qquad \text{når} \qquad (x,y,z) \in \partial W. 
    \end{align*}
    %
    Vis at $f(x,y,z) = 0$ for alle $(x,y,z) \in W$. 
\end{enumerate}

\oppgave[K2015, Oppgave 7]

La $S$ være delen av flata $z = 1 + x^2 + y^2$ som ligger inne i sylinderen $x^2
+ y^2 = 1$, og la $T$ være legemet begrenset av $S$ og planet $z = 2$.

\begin{enumerate}
  \item Gitt vektorfeltet $\vek{F} = f(x,y) \I + x g(x,y) \J + z \K$, der $f$ og
    $g$ har kontinuerlige partiellderiverte som oppfyller
    %
    \begin{equation*}
      \diffp{f}{x} + x \diffp{g}{y} = 0
    \end{equation*}
    %
    for alle $x, y$.

    Finn flateintegralet $\iint_S \vek{F} \cdot \dSS$, der enhetsnormalvektoren
    til $S$ har \emph{negativ} $\K$ komponent.
\end{enumerate}

\oppgave[V2014, Oppgave 7]

La $W$ være området i $\R^3$ gitt ved
%
\begin{equation*}
  W := \{ (x,y,z) \in \R^3 \mid 0 \leq x \leq y \ \text{og} \ x^2 + y^2 \leq z \leq 1 \}\,
\end{equation*}
%
med rand $S$ orientert slik at normalvektorene peker ut av $W$. La
$\vek{F}$ være vektorfeltet gitt ved $\vek{F}(x,y,z) = (y,x^2,xz)$. Bestem
%
\begin{equation*}
  \iint_{S} \vek{F} \cdot \dSS
\end{equation*}

\oppgave[K2013, Oppgave 7] La $T$ være legemet gitt ved $x \geq 0,$, $y \geq 0$
og $0 \leq z \leq 1 - \sqrt{x^2 + y^2}$. La $S$ være overflaten til $T$
orientert med enhetsnormal som peker ut av legemet $T$.

\begin{enumerate}
    \stepcounter{enumi}
  \item Finn $\displaystyle (x + \cos y, y + \sin z, z + e^x) \cdot \dSS$
\end{enumerate}

\oppgave[K2012, Oppgave 6] Legemet $T$ er avgrenset av paraboloiden $z = 4x^2 +
4y^2$ og planet $z = 4$. La $S$ være den delen av overflaten til $T$ som ligger
på paraboloiden $z = 4x^2 + 4y^2$, og la $\vek{n}$ være enhetsnormalen til $S$
som peker ut av legemet $T$. La videre $\vek{F}$ være vektorfeltet definert ved
%
\begin{equation*}
  \vek{F}(x,y,z)
  = \frac{yz}{8\pi} \I - \frac{x}{2\pi}\J + \frac{z}{4}\K
\end{equation*}
%
\begin{enumerate}
    \stepcounter{enumi}
  \item Regn ut $\displaystyle \iint_S \vek{F} \cdot \vek{n} \dS$.
\end{enumerate}


\oppgave[V2012, Oppgave 5] Finn fluksen av vektorfeltet $\vek{F}(x,y,z) = (\sin
y)^2\I + y\J + xz\K$ inn i enhetskula med sentrum i origo. 


%%% Local Variables:
%%% mode: latex
%%% TeX-master: "main"
%%% End:


\seksjon{Stokes' teorem}


\oppgave[V2017, Oppgave 5]

\begin{enumerate}
    \stepcounter{enumi}
  \item La $S$ være ei kuleflate i $\R^3$ og $\vek{F} \colon \R^3 \to \R^3$ ett
    glatt vektorfelt. Vis at
    %
    \begin{equation*}
     \iint_S \curl \vek{F} \cdot \dSS = 0.
    \end{equation*}
\end{enumerate}


\oppgave[V2015, Oppgave 9]

La $S$ være sylinderflata $\{(x,y,z) \in \R^3 \mid x^2 + y^2 = 1 \text{og} \ 0
\leq z \leq 1\}$. Finn
%
\begin{equation*}
  \iint_S \curl \vek{v} \cdot \dSS
\end{equation*}
%
når $\vek{v} = (v_1, v_2, v_3)$ er et $C^2$-vektofelt i $\R^3$ der  $v_1$ og
$v_2$ ikke avhenger av $z$.


\oppgave[K2014]

I denne oppgaven studerer vi de to flatene
%
\begin{align*}
  S_1 & := \{ (x,y,z) \in \R^3 \mid x^2 + y^2 + z^2 = 1, \ z \geq 1\} \\
  S_2 & := \{ (x,y,z) \in \R^3 \mid x^2 + y^2 + \frac{z^2}{4} = 1, \ z \geq 1\} 
\end{align*}
%
som berre er orientert slik at normalvektoren til $S_1$ i $(0,0,1)$ og $S_2$ i
$(0,0,2)$ er $\K$.

\begin{enumerate}
  \item Begrun at hvis $\vek{F}$ er et $C^1$ vektorfelt i $\R^3$, så er
    %
    \begin{equation*}
      \iint_{S_1} \curl \vek{F} \cdot \dSS = \iint_{S_2} \curl \vek{F} \cdot \dSS\,.
    \end{equation*}
    %
  \item Finn verdien av integralene når $\vek{F}(x,y,z) = -2y \I + x\J + e^{x^2+z}\K$.
\end{enumerate}

\oppgave[V2013, Oppgave 6] I denne oppgaven er

\begin{itemize}
  \item $S$ sfæren (kuleoverflaten) med sentrum i $(0,0,3)$ og radius $5$,
    orientert slik at normalvektoren peker ut av kula.
  \item $S^+$ den delen av $S$ som ligger over $xy$-planet, med samme
    orientering som $S$.
  \item $C$ er randen til $S^+$, med positiv orientering i samsvar med
    orienteringen til $S^+$.
  \item $\vek{F}(x,y,z) = (16 - x^2 - y^2)\I + z\J + (x + y + z)\K$.
\end{itemize}

\begin{enumerate}
  \item Finn en parametrisering for kurven $C$ (husk orienteringen). Lag en
    skisse som viser $S^+$, $C$ og $T$. Merk på orienteringene til $S^+$ og $C$.
    \stepcounter{enumi}
  \item Finn fluksen til $\curl \vek{F}$ gjennom flaten $S^+$. Altså beregn
    %
    $ \displaystyle 
    \int_{S^+} \curl \vek{F} \cdot \dSS.
    $
\end{enumerate}

\oppgave[K2012, Oppgave 6] Legemet $T$ er avgrenset av paraboloiden $z = 4x^2 +
4y^2$ og planet $z = 4$. La $S$ være den delen av overflaten til $T$ som ligger
på paraboloiden $z = 4x^2 + 4y^2$, og la $\vek{n}$ være enhetsnormalen til $S$
som peker ut av legemet $T$. La videre $\vek{F}(x,y,z)$ være vektorfeltet definert ved
%
\begin{equation*}
  \vek{F}(x,y,z)
  = \frac{yz}{8\pi} \I - \frac{x}{2\pi}\J + \frac{z}{4}\K.
\end{equation*}
%
\begin{enumerate}
    \setcounter{enumi}{2}
  \item Regn ut $\displaystyle \iint_S \curl \vek{F} \cdot \vek{n} \dS$.
\end{enumerate}

\oppgave[V2012, Oppgave 5] La $S$ være halvkula $x^2 + y^2 + z^2 = 4$ med $z\leq
0$ orientert nedover. La $\vek{F}(x,y,z) = x \tan (z/4) \I + x e^{e^{z^4}}\J +
xyz \K$. Finn integralet $\displaystyle \iint_S \curl \vek{F} \cdot \dSS$.



%%% Local Variables:
%%% mode: latex
%%% TeX-master: "main"
%%% End:


\begin{tikzpicture}
  \foreach \i in {1,...,\arabic{section}}
  {
    \node at (0, -\i) {\nameref{sec:\i}: \arabic{problem\i}};
  };
\end{tikzpicture}

\begin{figure}[htbp!]
  \centering
    \begin{tikzpicture}
      \begin{axis}[ xbar, y=-0.5cm, bar width=0.3cm, y axis line style = {
          opacity = 0 }, axis x line = none, tickwidth = 0pt, enlarge y limits =
        0.2, enlarge x limits = 0.02, ytick = data, nodes near coords, symbolic
        y coords = {%
          Grenseverdier,
          Gradient,
          Epsilon-Delta,
          Kjerneregelen,
          Taylor-approksimasjon,
          Tangenter,
          Kritiske punkter,
          Optimering,
          Linjeintegral,
          Konservative vektorfelt,
          Dobbelintegraler,
          Integrasjonsrekkefølge,
          Greens teorem,
          Trippelintegral,
          Divergensteoremet,
          Stokes teorem
          }
          ]
          \addplot[fill=black] coordinates {%
            (\arabic{problem1},Grenseverdier)
            (\arabic{problem2},Gradient)
            (\arabic{problem3},Epsilon-Delta)  
            (\arabic{problem4},Kjerneregelen)  
            (\arabic{problem5},Taylor-approksimasjon)  
            (\arabic{problem6},Tangenter)  
            (\arabic{problem7},Kritiske punkter)  
            (\arabic{problem8},Optimering)
            (\arabic{problem9},Linjeintegral)
            (\arabic{problem10},Konservative vektorfelt)
            (\arabic{problem11},Dobbelintegraler)  
            (\arabic{problem12},Integrasjonsrekkefølge)  
            (\arabic{problem13},Greens teorem)  
            (\arabic{problem14},Trippelintegral)  
            (\arabic{problem15},Divergensteoremet)  
            (\arabic{problem16},Stokes teorem)};
        % \legend{Topics, Posts}
      \end{axis}
    \end{tikzpicture}
  \end{figure}

\end{document}
